% P15-U1 / N-x-008: independence, product laws and moment consequences.
\chapter{Independent Events and Random Variables}
\label{sec:p15-independent-random-variables}

Independence is a factorization property of a specified probability model.
It concerns every required combination of measurable events, not just
uncorrelated numerical observations or a few observed pairs. This chapter
develops that property for events, random elements and sigma-algebras,
then proves when measurable transformations, grouped variables and product
calculations preserve it.

The principal source is Durrett's author Version 5 dated 11 January 2019,
Section 2.1 through Section 2.1.2, including the definitions and results
numbered 2.1.1--2.1.14 \cite[pp.~43--49]{DurrettPTEAuthor2019}.
The source's finite construction, countable-valued criterion, generator
counterexample and dyadic coin construction are also supplied
\cite[p.~53; exercises~2.1.2, 2.1.6, 2.1.9, 2.1.15]{DurrettPTEAuthor2019}.
The final moment and mean-square results include results 2.2.1--2.2.3
\cite[pp.~56--57]{DurrettPTEAuthor2019}.
The event, integration and product prerequisites are the actual
constructions in Chapters~\ref{sec:p15-measurable-probability},
\ref{sec:p15-integration-expectation} and~\ref{sec:p15-product-measures}.

\section{Families and their finite subfamilies}
\label{sec:p15-independent-families}

\begin{hbdefinition}{def:p15-general-independence}{Independence of events, collections and sigma-algebras}
Fix a probability space \((\Omega,\mathcal F,\mathbb P)\), and let
\(I\) be a finite or countably infinite index set. A family of events
\((A_i)_{i\in I}\) is \emph{mutually independent}, or simply
\emph{independent}, if for every finite \(J\subseteq I\),
\[
 \mathbb P\left(\bigcap_{i\in J}A_i\right)
                         =\prod_{i\in J}\mathbb P(A_i).
\]
The empty intersection is \(\Omega\), and the empty product is one.
The family is \emph{pairwise independent} when the equality is required
only for the two-element subfamilies. Even for a finite family, checking
only the intersection of all its events is not this definition.

More generally, collections \(\mathcal C_i\subseteq\mathcal F\)
are independent if for every finite \(J\subseteq I\) and every
selection \(A_i\in\mathcal C_i\), \(i\in J\), the same factorization
holds. In particular this defines independent sub-sigma-algebras
\(\mathcal F_i\subseteq\mathcal F\). In a finite family of
sigma-algebras it suffices to test one event from every sigma-algebra:
setting the unused events equal to \(\Omega\) recovers each subfamily.

Independence of an infinite family means independence of every finite
subfamily. This definition does not itself construct a measure on a
space of infinite sequences, and it does not replace finite products
by a formal infinite product of event probabilities.
\index{independence!mutual}\index{independence!sigma-algebras}
\end{hbdefinition}

\begin{hbdefinition}{def:p15-independent-random-elements}{Independence of measurable observations}
For measurable spaces \((S_i,\mathcal S_i)\), let
\(X_i:\Omega\to S_i\) be random elements. They are independent
if, for every finite \(J\subseteq I\) and every \(B_i\in\mathcal S_i\),
\[
 \mathbb P\left(\bigcap_{i\in J}\{X_i\in B_i\}\right)
                      =\prod_{i\in J}\mathbb P(X_i\in B_i).
\]
For real random variables, the value sigma-algebra is
\(\mathcal B(\mathbb R)\), as in the source. Write
\(\sigma(X_i)=\{X_i^{-1}(B):B\in\mathcal S_i\}\), the already
proved sigma-algebra generated by that observation. Pairwise independence
again tests only pairs. \emph{Identically distributed} means that all
variables have the same measurable value space and the same law;
\emph{independent and identically distributed}, abbreviated \emph{i.i.d.},
requires both properties. Equal laws alone do not imply independence.
\index{independence!random elements}\index{i.i.d.}
\end{hbdefinition}

\begin{hbtheorem}{thm:p15-independence-basic}{proved=annex:p15.independence-basic}{Generated information, complements and indicators}
\par\nobreak
For the preceding finite or countable families:
\begin{enumerate}
\item The \(X_i\) are independent exactly when the sigma-algebras
      \(\sigma(X_i)\) are independent. If independent sigma-algebras
      \(\mathcal F_i\) make \(X_i\) measurable, then the \(X_i\)
      are independent. Taking subcollections of indices or smaller
      sigma-algebras preserves independence.
\item Complementing any selection of independent events preserves
      independence. Events \(A_i\) are independent exactly when their
      indicators \(\mathbf1_{A_i}\) are independent, and exactly when
      their sigma-algebras \(\sigma(\{A_i\})\) are independent.
      In particular all four choices of complement for an independent
      pair remain independent.
\item An event is independent of itself exactly when its probability
      is zero or one. Every such event is independent of every event.
      If \(\mathbb P(B)>0\), the two-event definition is equivalent
      to \(\mathbb P(A\mid B)=\mathbb P(A)\), where the conditional
      probability is \(\mathbb P(A\cap B)/\mathbb P(B)\).
\item Replacing each independent real \(X_i\) by a real random variable \(Y_i\)
      equal to it a.s. preserves independence and each marginal law.
      The replacement must be measurable, and each equality has a
      measurable null exceptional set; no completion of \(\mathcal F\)
      is assumed.
\end{enumerate}
These definitions agree with the earlier finite probability definitions
on their common domains.
\end{hbtheorem}

\section{Checking independence on generating events}
\label{sec:p15-independence-generators}

\begin{hbtheorem}{thm:p15-independence-generators}{proved=annex:p15.independence-generators}{Independent pi-systems generate independent sigma-algebras}
Adjoining \(\Omega\) to each collection \(\mathcal C_i\)
preserves independence. For a finite family containing \(\Omega\)
in every collection, the complete-family factorization is equivalent
to independence of all its subfamilies.

Suppose now that the collections \(\mathcal C_i\), \(i\in I\),
are independent and each is a pi-system, meaning it is closed under
binary intersections. Then \((\sigma(\mathcal C_i))_{i\in I}\)
are independent. The finite theorem is Durrett's Theorem 2.1.7;
the countable-family conclusion follows by applying it to every finite
subfamily. The proof uses the full pi-lambda theorem already proved
in Theorem~\ref{thm:p15-measure-uniqueness}.
\end{hbtheorem}

\begin{hbtheorem}{thm:p15-independence-thresholds}{proved=annex:p15.independence-thresholds}{A distribution-function criterion}
Real random variables \(X_1,\ldots,X_n\), \(n\geq1\), are
independent exactly when
\[
 \mathbb P(X_1\leq t_1,\ldots,X_n\leq t_n)
                        =\prod_{i=1}^n\mathbb P(X_i\leq t_i)
\]
for all \(t_i\in\mathbb R\cup\{\infty\}\), where
\(\{X_i\leq\infty\}=\Omega\). It also suffices to check
\(t_i\in\mathbb Q\cup\{\infty\}\). The infinite threshold
permits a coordinate to impose no event restriction. For a countable
family, apply this test to each finite list of distinct indices.
\end{hbtheorem}

\begin{hbtheorem}{thm:p15-independent-blocks}{proved=annex:p15.independent-blocks}{Disjoint groups and measurable functions}
Let independent sigma-algebras \(\mathcal F_i\), \(i\in I\),
be given, and let \((I_j)_{j\in J}\) be pairwise disjoint subsets
of \(I\), with \(J\) finite or countable. Put
\[
 \mathcal G_j\defeq\sigma\left(\bigcup_{i\in I_j}\mathcal F_i\right).
\]
Then the sigma-algebras \(\mathcal G_j\) are independent.
An empty group gives the trivial sigma-algebra \(\{\varnothing,\Omega\}\).
Any family of random elements \(Y_j\) such that \(Y_j\) is
\(\mathcal G_j\)-measurable is consequently independent.

In particular, measurable functions \(f_i(X_i)\) of independent
random elements are independent. For positive integers \(r,m(1),\ldots,m(r)\),
mutually independent real variables
\(X_{i,j}\), \(1\leq i\leq r\), \(1\leq j\leq m(i)\),
and Borel functions \(f_i:\mathbb R^{m(i)}\to\mathbb R\),
the variables \(f_i(X_{i,1},\ldots,X_{i,m(i)})\) are independent.
For countably infinite groups, a proposed output's measurability with
respect to its indicated \(\mathcal G_j\) is still required.
\end{hbtheorem}

\begin{hbexample}{def:p15-independent-past-future}{A future increment and an earlier crossing event}
Let \(X_1,\ldots,X_n\) be mutually independent real variables,
\(S_k=\sum_{i=1}^kX_i\), and \(1\leq m<n\). For real \(a\),
\[
 S_n-S_m=\sum_{i=m+1}^nX_i
 \quad\hbox{and}\quad
 \mathbf1_{\{\max_{1\leq k\leq m}S_k>a\}}
\]
are independent. The increment is a Borel function of the future block
\(X_{m+1},\ldots,X_n\). The indicator is a Borel function of
\(X_1,\ldots,X_m\), since finite sums and maxima are continuous
and the strict threshold defines a Borel event. The blocks are disjoint,
so the preceding theorem applies. No infinite sum or stopping-time theorem
is used in this conclusion.
\end{hbexample}

\section{Product laws and expectation}
\label{sec:p15-independence-laws-expectation}

\begin{hbtheorem}{thm:p15-independent-product-law}{proved=annex:p15.independent-product-law}{Joint laws, finite construction and countable value sets}
\par\nobreak
Let \(n\geq1\) and let \(X_i:\Omega\to(S_i,\mathcal S_i)\), \(1\leq i\leq n\),
be random elements with laws \(\mu_i\). They are independent exactly
when their joint law on the finite product sigma-algebra is
\(\mu_1\times\cdots\times\mu_n\).

Conversely, for any finite list of probability spaces
\((S_i,\mathcal S_i,\mu_i)\), the coordinate maps on their constructed
product probability space are independent and have the specified laws.
For prescribed admissible real distribution functions \(F_i\), take
the constructed Stieltjes probability measures \(\mu_{F_i}\).
This gives real independent variables whose joint probabilities on
half-open rectangles with real endpoints \(a_i<b_i\) are
\[
 \mathbb P(a_i<X_i\leq b_i\text{ for all }i)
                       =\prod_{i=1}^n(F_i(b_i)-F_i(a_i)).
\]

If \(X_i\) take their values in nonempty countable sets \(S_i\)
with all subsets measurable, it suffices to check
\[
 \mathbb P(X_1=x_1,\ldots,X_n=x_n)
                     =\prod_{i=1}^n\mathbb P(X_i=x_i)
               \qquad(x_i\in S_i).
\]
This recovers the earlier finite-valued criterion. All constructions in
this theorem concern finitely many coordinates.
\end{hbtheorem}

\begin{hbtheorem}{thm:p15-independent-expectations}{proved=annex:p15.independent-expectations}{Integration with product laws and factorization of moments}
Let independent real variables \(X,Y\) have laws \(\mu,\nu\).
For a Borel function \(h:\mathbb R^2\to[0,\infty]\), or a real
Borel function with \(\mathbb E[|h(X,Y)|]<\infty\),
\[
 \mathbb E[h(X,Y)]
      =\int_{\mathbb R^2}h\,d(\mu\times\nu)
      =\int_{\mathbb R}\left(\int_{\mathbb R}h(x,y)\,d\mu(x)\right)d\nu(y).
\]
The reversed iterated order has the same value. In the signed case,
inner integrals are assigned a measurable finite value on their null
exceptional sets as in Theorem~\ref{thm:p15-fubini}.

For Borel real functions \(f,g\), if both are nonnegative, or if
\(\mathbb E|f(X)|<\infty\) and \(\mathbb E|g(Y)|<\infty\), then
\[
 \mathbb E[f(X)g(Y)]=\mathbb E[f(X)]\mathbb E[g(Y)].
\]
The signed case includes absolute integrability of the product.
For \(n\geq1\) independent real \(X_1,\ldots,X_n\), either all nonnegative
or all individually integrable,
\[
 \mathbb E\left[\prod_{i=1}^nX_i\right]
                         =\prod_{i=1}^n\mathbb E[X_i].
\]
In the integrable case the product is integrable and
\(\mathbb E\prod_i|X_i|=\prod_i\mathbb E|X_i|<\infty\).
Nonnegative equalities permit infinity and use \(0\cdot\infty=0\).

Conversely, factorization \(\mathbb E[f(X)g(Y)]=
\mathbb E[f(X)]\mathbb E[g(Y)]\) for every pair of bounded Borel
real \(f,g\) implies independence. The corresponding finite-product
test for bounded functions on every finite subfamily characterizes
independence of a countable real family.
\end{hbtheorem}

\begin{hbdefinition}{def:p15-covariance-correlation}{Covariance, correlation and uncorrelated variables}
For real \(X,Y\) with finite second moments, define
\[
 \operatorname{Cov}(X,Y)\defeq
       \mathbb E[(X-\mathbb E X)(Y-\mathbb E Y)].
\]
The following theorem verifies existence of this integral. The variables
are \emph{uncorrelated} if this covariance is zero. When both variances
are strictly positive, their \emph{correlation} is
\[
 \operatorname{Corr}(X,Y)\defeq
 \frac{\operatorname{Cov}(X,Y)}
      {\sqrt{\operatorname{Var}(X)\operatorname{Var}(Y)}}.
\]
It is undefined by this formula if either variance is zero. A family
with finite second moments is uncorrelated when each distinct pair is.
\index{covariance}\index{correlation}\index{uncorrelated variables}
\end{hbdefinition}

\begin{hbtheorem}{thm:p15-independent-moments}{proved=annex:p15.independent-moments}{Covariance and the variance of a finite sum}
For finite-second-moment real variables, covariance exists and
\[
 \operatorname{Cov}(X,Y)=\mathbb E[XY]-\mathbb E[X]\mathbb E[Y],
 \qquad
 |\operatorname{Cov}(X,Y)|
       \leq\sqrt{\operatorname{Var}(X)\operatorname{Var}(Y)}.
\]
Thus defined correlations lie in \([-1,1]\). Independence implies
zero covariance under these moment hypotheses.
For \(X_1,\ldots,X_n\) with finite second moments and real constants
\(a_1,\ldots,a_n\), the sum has finite second moment and
\[
 \operatorname{Var}\left(\sum_{i=1}^n a_iX_i\right)
  =\sum_{i=1}^n a_i^2\operatorname{Var}(X_i)
      +2\sum_{1\leq i<j\leq n}a_ia_j\operatorname{Cov}(X_i,X_j).
\]
For uncorrelated variables the second sum is zero. Pairwise independence
is sufficient for this conclusion; mutual independence is unnecessary.
\end{hbtheorem}

\section{What weaker checks do not establish}
\label{sec:p15-independence-boundaries}

\begin{hbtheorem}{thm:p15-independence-counterexamples}{proved=annex:p15.independence-counterexamples}{Pairwise independence, generator closure and zero correlation}
\par\nobreak
The following explicitly constructed finite models give four boundaries.
\begin{enumerate}
\item On \(\{0,1\}^3\) with mass \(1/8\) at every outcome,
      let \(X_i\) be the coordinates and set
      \(A_1=\{X_2=X_3\}\), \(A_2=\{X_3=X_1\}\),
      \(A_3=\{X_1=X_2\}\). The coordinates are independent.
      Each \(A_i\) has probability \(1/2\); each pairwise
      intersection has probability \(1/4\), but the triple intersection
      also has probability \(1/4\), instead of \(1/8\).
      Thus the events, and likewise their indicators, are pairwise
      independent but not mutually independent. This is the source's
      Example 2.1.4 with its underlying model supplied.
\item On \(\Omega=\{1,2,3,4\}\) with uniform mass, take
      \(A=\{1,2\}\), \(B=\{1,3\}\), \(C=\{1,4\}\).
      The two collections \(\{A,B\}\) and \(\{C\}\) are
      independent, but their generated sigma-algebras are not.
      The first collection fails the pi-system requirement.
\item Let the joint masses of real \(X,Y\) be
      \[
       \begin{array}{c|ccc}
          &Y=1&Y=0&Y=-1\\\hline
        X=1&0&a&0\\
        X=0&b&c&b\\
        X=-1&0&a&0
       \end{array}
      \]
      with \(a,b>0\), \(c\geq0\), and \(2a+2b+c=1\).
      Then \(\mathbb E X=\mathbb E Y=\mathbb E[XY]=0\),
      \(\operatorname{Var}(X)=2a\), \(\operatorname{Var}(Y)=2b\),
      and the correlation is zero. Nevertheless
      \(\mathbb P(X=1,Y=1)=0<ab\).
      This retains the source's Example 2.1.14 and its full parameter range.
\item A \(\{0,1\}\)-valued random variable \(Z\) with
      \(\mathbb P(Z=1)=1/2\), called a fair Bernoulli variable,
      is not independent of itself. Reusing it in two output groups
      therefore does not satisfy the disjoint-group conclusion, even
      though each output is a measurable function of its own named group.
\end{enumerate}
\end{hbtheorem}

\begin{hbmethod}{def:p15-independence-method}{Establishing an independence calculation}
Specify the underlying probability measure and the family whose
independence is required. For finitely many real variables, establish
factorization on the generating threshold events, or prove that the
joint law is the product of its marginals. On countable value sets,
checking all joint point masses is sufficient. Use the disjoint-group
theorem before replacing original variables by block functions. To factor
a signed expected product, also establish the stated absolute moments.
Pairwise independence, or zero covariance for every distinct pair,
justifies variance addition when the second moments exist; these
conditions do not establish the full joint-law factorization for three
or more variables. An assertion about observations from data additionally
requires an inference model; these theorems concern the specified law.
\end{hbmethod}

\section{An actual independent sequence and a mean-square consequence}
\label{sec:p15-independent-sequence}

\begin{hbtheorem}{thm:p15-dyadic-independent-coins}{proved=annex:p15.dyadic-independent-coins}{Independent fair bits on the constructed unit interval}
On the constructed uniform probability space \((0,1)\), with
\(\mathbb P=\mathbb U\), define
\[
 Y_n(u)=\lfloor2^nu\rfloor-2\lfloor2^{n-1}u\rfloor,
                      \qquad n\geq1.
\]
Here \(\lfloor t\rfloor\) is the greatest integer at most \(t\).
These are measurable \(\{0,1\}\)-valued variables; equivalently,
\(Y_n(u)\) is one precisely when \(\lfloor2^nu\rfloor\) is odd,
as in the source's Exercise 2.1.15. The countable family is mutually
independent, with \(\mathbb P(Y_n=0)=\mathbb P(Y_n=1)=1/2\).
For any finite \(J\subseteq\mathbb N_+\) and any prescribed
\(\varepsilon_j\in\{0,1\}\),
\[
 \mathbb P(Y_j=\varepsilon_j\text{ for all }j\in J)=2^{-\#J}.
\]
The floor convention specifies every dyadic endpoint and every outcome
in \((0,1)\). This constructs a particular infinite independent
sequence using a measure already available in the manuscript.
\end{hbtheorem}

\begin{hbtheorem}{thm:p15-independent-l2-weak-law}{proved=annex:p15.independent-l2-weak-law}{Moment control and the mean-square weak law}
If real random variables \(Z_n\) satisfy
\(\mathbb E|Z_n|^p\to0\) for some real \(p>0\), then
\(Z_n\to0\) in probability.

Let real \(X_1,X_2,\ldots\) have finite second moments, be
pairwise uncorrelated, and satisfy
\(\mathbb E X_i=\mu\) and \(\operatorname{Var}(X_i)\leq C<\infty\)
for every \(i\). For \(S_n=\sum_{i=1}^nX_i\),
\[
 \mathbb E\left[\left(S_n/n-\mu\right)^2\right]\leq C/n,
 \qquad
 \mathbb P(|S_n/n-\mu|\geq\varepsilon)
                     \leq\frac{C}{n\varepsilon^2}\quad(\varepsilon>0).
\]
Consequently \(S_n/n\to\mu\) in mean square and in probability.
Mean-square convergence means that the displayed expected square tends
to zero. Applied to the constructed fair bits, \(\mu=1/2\),
\(C=1/4\), giving the bound \(1/(4n\varepsilon^2)\) for the
deviation of their empirical proportion from one half.
\end{hbtheorem}

\begin{proofboundary}{def:p15-independence-next}{Convolutions and the general infinite-product construction}
All results stated here have full annex proofs.
\Cref{sec:p15-convolution-gamma-normal} develops the source's
convolutions, densities and gamma and normal sum presentations in
results 2.1.15--2.1.20, with both gamma proof routes.
The countable extension in \cref{sec:p15-countable-product-measures}
supplies Theorem 2.1.21 and Appendix A.3, followed by the full
measurable-space representation of Theorem 2.1.22. It also proves that
the fair-bit construction above has the canonical infinite product law.
The extension's stated measurable-space hypotheses remain essential.
The source's additional independence exercises and later weak, strong
and zero-one laws also retain their separate required statements and proofs.
\end{proofboundary}
